% TeleCon firmware sketches — which .ino to flash.
% Build: pdflatex -output-directory=../build telecon_firmware_sketches.tex
% Delivered PDF (Android Codes asset): ../build/firmware_sketches_eng.pdf
\documentclass[12pt,a4paper]{article}

\usepackage[margin=2.4cm]{geometry}
\usepackage[T1]{fontenc}
\usepackage[utf8]{inputenc}
\usepackage{lmodern}
\usepackage{booktabs}
\usepackage{hyperref}
\usepackage{xcolor}
\usepackage{colortbl}
\usepackage{array}
\usepackage{enumitem}
\usepackage{needspace}
\usepackage{etoolbox}

\preto\section{\Needspace{8\baselineskip}}

\hypersetup{
  colorlinks=true,
  linkcolor=black,
  urlcolor=blue,
  pdftitle={TeleCon firmware sketches},
  pdfauthor={TeleCon}
}

\definecolor{TblHead}{HTML}{0F6F73}
\definecolor{TblChan}{HTML}{D5F5E3}
\definecolor{TblSw}{HTML}{FCF3CF}
\definecolor{TblTelem}{HTML}{D4E6F1}
\definecolor{TblKnob}{HTML}{F5CBA7}

\newcommand{\thc}[1]{\textcolor{white}{\bfseries #1}}
\newcommand{\chip}[2]{\colorbox{#1}{\strut\,\scriptsize #2\,}}
\newcommand{\ui}[1]{\textsf{\textbf{#1}}}

\title{TeleCon firmware sketches\\[0.3em]
\large Which \texttt{.ino} to flash}
\author{ESP32 DevKit over Bluetooth, and ESP32-CAM for video}
\date{Version 1.0 \quad 2026}

\begin{document}
\maketitle

\begin{abstract}
The Android app downloads one ZIP per connection from the firmware
repository. Each ZIP is one Arduino sketch, and the sketch is the
\texttt{.ino} with the same name. This note names every Bluetooth control
sketch and the ESP32-CAM sketch. The camera connects only with a DevKit
over Bluetooth and a separate ESP32-CAM for video.
\end{abstract}

{\sffamily\scriptsize
\chip{TblChan}{DevKit + Bluetooth}
\chip{TblTelem}{CAM video only + DevKit}}

\tableofcontents
\vspace{1em}

\section{The only way to connect the camera}

The picture does not come from the same board as the sticks. Two boards
are required, and the app has no other camera layout.

\begin{center}
\renewcommand{\arraystretch}{1.28}
\setlength{\tabcolsep}{5pt}
\begin{tabular}{>{\raggedright\arraybackslash}p{3.6cm}
                >{\raggedright\arraybackslash}p{5.4cm}
                >{\raggedright\arraybackslash}p{5.2cm}}
\rowcolor{TblHead}
\thc{Role} & \thc{Board} & \thc{Sketch} \\
\rowcolor{TblChan}
Sticks & ESP32 DevKit & One of the three Bluetooth \texttt{.ino} files \\
\rowcolor{TblTelem}
Picture & ESP32-CAM & \texttt{TeleCon\_CAM\_\-SoftAP\_Video.ino} \\
\end{tabular}
\end{center}

Flash the video sketch on the CAM and one Bluetooth sketch on the DevKit.
On the phone, keep the Bluetooth mode pointed at the DevKit and turn on
the camera overlay. The phone joins Wi-Fi \texttt{TeleCon-RC-CAM-Starter},
password \texttt{telecon1234}, for the picture only. Sticks stay on
Bluetooth. One ESP32-CAM does not carry them, and control does not travel
on the camera Wi-Fi.

With no camera, the DevKit and one of the three Bluetooth sketches are
enough. That board has no video.

\section{Bluetooth control on an ESP32 DevKit}

Board in the app: \ui{DevKit}. No camera on this board. Open the sketch
folder in the Arduino IDE (the folder that contains the \texttt{.ino}) and
upload to an ESP32 Dev Module.

\begin{center}
\renewcommand{\arraystretch}{1.28}
\setlength{\tabcolsep}{4pt}
\begin{tabular}{>{\raggedright\arraybackslash}p{4.3cm}
                >{\raggedright\arraybackslash}p{4.6cm}
                >{\raggedright\arraybackslash}p{5.3cm}}
\rowcolor{TblHead}
\thc{ZIP / \texttt{.ino}} & \thc{Phone setting} & \thc{On the air} \\
\rowcolor{TblChan}
\texttt{TeleCon\_Classic\_Simple} & Classic + Simple & Name \texttt{ESP32-TC-RC-BT-Simple}, PIN \texttt{1234}. Handshake \texttt{proto,simple}. Text lines \texttt{RC:CTRL}, \texttt{RC:BTN}, \texttt{RC:DATA}, \texttt{RC:PLOT}. \\
\rowcolor{TblChan}
\texttt{TeleCon\_Classic\_Binary} & Classic + Binary & Name \texttt{ESP32-TC-RC-BT-Binary}. Handshake \texttt{proto,binary}. Frames \texttt{AA 55}, \texttt{BB 66}, \texttt{CC 11/22/33}. \\
\rowcolor{TblChan}
\texttt{TeleCon\_BLE\_Binary} & BLE Binary & Name \texttt{ESP32-TC-RC-BLE-Binary} (Nordic UART). Same binary frames as Classic Binary. \\
\end{tabular}
\end{center}

\paragraph{\texttt{TeleCon\_Classic\_Simple.ino}}
Single file for a first board. Put motor and servo code in
\texttt{onControl()} and \texttt{onButton()} where the file says
\texttt{YOUR CODE HERE}. A steering servo on GPIO 13 is already wired in
the sketch: stick \texttt{rx} $= 0$ is a $1500$\,\textmu s pulse. Eight
telemetry bytes (phone CH1--CH8) are chosen in
\texttt{kTelemetryChannels[]} in this same file. With \texttt{ENABLE\_PLOT}
at 1 the board sends \texttt{RC:PLOT} about every 50\,ms.
The code guide for this sketch is
\texttt{telecon\_classic\_simple\_code.tex} in this documentation set.

\paragraph{\texttt{TeleCon\_Classic\_Binary.ino}}
The \texttt{.ino} only starts Bluetooth and the vehicle modules. Pins,
PWM, and the eight telemetry channels live in the other files in that
folder. Edit \texttt{pins.h} for your driver and
\texttt{TelemetryChannels.h} for what CH1--CH8 report. The phone setting
must be Classic + Binary before connect. The wire format is the control
protocol document, not the text lines of the Simple sketch.

\paragraph{\texttt{TeleCon\_BLE\_Binary.ino}}
Same vehicle modules and the same \texttt{TelemetryChannels.h} table as
Classic Binary. The link is BLE instead of Classic SPP. Choose BLE Binary
in the app. Do not pair it as a Classic device.

These three sketches are also the control half of the two-board camera
layout below. Flash only one of them on the DevKit.

\section{ESP32-CAM, video only}

Use this when a DevKit already runs one of the Bluetooth sketches and you
want a picture from a second board.

\paragraph{\texttt{TeleCon\_CAM\_SoftAP\_Video.ino}}
ZIP \texttt{TeleCon\_CAM\_SoftAP\_Video.zip}. Flash an ESP32-CAM
(AI-Thinker). This sketch serves video and does not accept sticks.

\begin{itemize}[leftmargin=1.4cm]
  \item SoftAP \texttt{TeleCon-RC-CAM-Starter}, password \texttt{telecon1234},
        address \texttt{192.168.4.1}.
  \item Picture: \texttt{http://192.168.4.1/stream} (MJPEG) or
        \texttt{/capture} (one JPEG).
  \item Stream quality: \texttt{http://192.168.4.1/camconfig}.
  \item No control port. There is no \texttt{RC:CONNECT} on this board.
\end{itemize}

On the phone, turn on the camera overlay and keep the Bluetooth connection
pointed at the DevKit (Classic Simple, Classic Binary, or BLE Binary).
The phone uses Wi-Fi for the CAM and Bluetooth for the sticks at the same
time. Do not put stick control on the ESP32-CAM.

\section{What the ZIP contains}

The app downloads \texttt{zips/TeleCon\_<name>.zip}. Unzip it and open the
folder that contains the \texttt{.ino} of the same name. Arduino compiles
every \texttt{.cpp} and \texttt{.h} in that folder, so a binary sketch is
the \texttt{.ino} plus those modules. Do not copy \texttt{SteerCenter}
back in: stock sketches do not store a steering center. Trim stays on the
phone, and \texttt{rx} $= 0$ is $1500$\,\textmu s.

\end{document}
